\chapter{Related Work} \label{ch:related}

\parab{On-demand Streaming.} A line of work has explored on-demand volumetric video streaming~\cite{vivo, groot, vues, yuzu, hermes, patchvvc}, and focused on reducing data rate and improving decoding efficiency. 
%
ViVo~\cite{vivo} culls occluded points or those outside the predicted frustum.
% 
Groot~\cite{groot} employs parallelism to improve Draco decoding efficiency and achieve 30~fps. 
%
Fumos~\cite{fumos} improves data rate by exploiting inter-frame redundancy using a neural codec and supports high frame rate decoding. 
%
In contrast, \livo avoids point cloud compression to achieve performant two-way live streaming.
%On-demand streaming relies on heavy pre-processing such as constructing index structure like octree/kd-tree, encoding point cloud at different density levels, occlusion-aware indexing, inter-frame redundancy, \etc, that are computationally expensive to be applied to live streaming.

\parab{Live Streaming and Conferencing.}
%There are a few live volumetric streaming systems compared to on-demand streaming. 
%
\cref{tab:comparison} explains how \livo differs from related work. We discuss salient differences below.
%
LiveScan3D~\cite{livescan3d} uses a fast binary compression method, zstd, which achieves low compression ratios.
%
%One-way live streaming does not have as stringent end-to-end latency requirements as two-way live streaming.
%
%is one of the earliest works to capture a scene using multiple Kinect cameras and stream live to the receiver. It can only support 15~FPS streaming since it uses ZSTD, which is a fast binary compression method, but achieves very low compression ratio for point cloud. 
%
Holoportation~\cite{holoportation} achieves conferencing by compressing a 3D representation of a single person using LZ4 binary encoding, which doesn't scale to larger scenes.
%
Starline~\cite{starline}, a conferencing system, streams upper body views in 3D over fixed quality 2D video streams; \cref{s:evaluation} describes how it differs from \livo.
%
FarfetchFusion~\cite{farfetchfusion} streams only the face and is optimized for mobile devices.
%
MeshReduce~\cite{meshreduce} supports full-scene conferencing using a mesh, but achieves a lower frame rate and lower quality (\cref{s:evaluation}).
%
MetaStream~\cite{metastream} uses Draco encoding to stream a single person at 30~fps by reducing the number of points in the capture point cloud. 
%
Unlike \livo, none of the live volumetric streaming systems support full-scene 2-way live streaming with bandwidth adaptation at 30~fps.
%To the best of our knowledge, \livo is the first to achieve this.

\parab{Other 3D Formats and Compression Standards.}
%
Besides point clouds, prior work attempts to capture 3D scenes as textured meshes~\cite{kinectfusion,fvv,holoportation,fusion4d,dynamic-fusion,meshreduce}.
%
%\agr{Maybe worth mentioning neural representations Nerf/Splats and just say they look amazing, but are currently too slow to train for interactive scenarios.  You could also hint that point cloud streaming could potentially be adapted to do splat streaming some day}
%
Meshes are generally higher quality representations that are much more bandwidth-intensive; future work can explore mesh-based two-way bandwidth-adaptive live streaming.
%
Besides Draco, prior work has proposed other point cloud compression techniques: G-PCC~\cite{gpcc}, V-PCC~\cite{vpcc}, and PCL~\cite{pcl}. 
%
Draco achieves faster compression with higher compression ratios than these, so we base our evaluations on it.
%
Recently, Hermes~\cite{hermes}, patchVVC~\cite{patchvvc}, and DeformStream~\cite{deformstream} have attempted to exploit inter-frame redundancy in volumetric videos, but these have high encoding latency, so can only be used for on-demand streaming.

\parab{Learned 3D Representations.}
Recent work uses learned 3D representations such as Neural Radiance Fields (NeRFs)~\cite{nerf} and Gaussian Splatting (\gs)~\cite{gsplat} for streaming and conferencing.
%
% Unlike point cloud and mesh, NeRF learns a neural network on a bunch of 2D images of a 3D scene and generates novel views in 2D from any perspective of the viewer's camera. 
% %
% Gsplat, an enhancement over NeRF, learns a bunch of 3D gaussians obtained from point clouds of the scene, and uses the gaussians' color and opacity to generate novel views. 
%
Research in this direction~\cite{nerf-nossdav, nerf-ems1, mga, swings, telepresence-tvcg, 3dgstream, queen, streamrf} have explored learned 3D representations for on-demand streaming. 
%
Tele-Aloha~\cite{tele-aloha} has used Gsplat for conferencing but constrains only to upper torso of a person like Starline~\cite{starline}.
%
Though these representations can generate high-quality rendering superior to point clouds or meshes, they suffer from high training and rendering overheads, making them unsuitable for live or conferencing systems. 
%
% We see promise in using learned representations, particularly Gsplat which is tightly linked with underlying point cloud. 
% %
% Many of the techniques used in \livo such as fast culling, depth encoding, and bandwidth splitting could potentially be adapted for Gsplat-based streaming in future.
%
%% \dgs requires high storage/bandwidth to store/stream which have been coped with by either in-training compression of gaussians or post-training compression of pre-trained models.

\parab{During-training \gs compression.}
Recent research has focused on integrating compression directly into an optimization loop within training to enhance model sparsity.
%
Compact3DGS~\cite{compact3dgs} identifies prunes redundant Gaussians (as do~\cite{recongs2025, scaffoldgs}) and replaces spherical harmonics with hash-based grids.
%
CompGS~\cite{compgs} utilizes quantization-aware training and opacity regularization.
%
4D-GS~\cite{4dgs} trains an MLP to predict motion across frames, while QUEEN~\cite{queen} learns quantized attribute residuals and uses a gating module to sparsify positions. 
%
%Others~\cite{recongs2025, scaffoldgs} use anchor-based structures to prune redundant Gaussian primitives to train a more compact model.
%
LapisGS~\cite{shi2025lapisgs} trains \gs{} in a layered progressive representation.
%
Relative to \gsnfs{}, these methods incur significantly higher overhead for compression.

\parab{Post-training \gs compression.}
% Post-training compression schemes aim to compress pre-trained models efficiently without the overhead of full retraining or finetuning of trained \gs. 
%
MesonGS~\cite{mesongs} prunes insignificant Gaussians and then applies octree and RAHT.
%
Some schemes~\cite{gauspcgc, entropygs} use learning-based methods to find the compression parameters for position and attributes.
%
Others~\cite{LTS, gpcc-adaptive} use point cloud compression techniques to compress trained \gs{}.
%
% But most post-training methods that compress in 3D suffer from significantly high encoding and decoding latency.
%
To address this, \vthree{}~\cite{v3volumetric} encodes \dgs{} in multiple 2D videos.
%
\gsnfs{} encodes and decodes much faster than these techniques.
%
%Though faster than methods that compress \gs in 3D, this approach suffers from higher size or lower quality, particularly for moderate to large scenes.

\parab{Quality Metrics.} 
%
Other work has explored different 3D objective quality metrics, point2point-PSNR~\cite{dong-tian}, PCQM~\cite{pcqm}, GraphSIM~\cite{graphsim}.
%
These compare point positions, and attributes such as colors, normals, and curvatures in 3D. 
%
We choose PointSSIM for measuring quality since it can measure both geometry and color distortions by directly extending the popular SSIM metric to 3D. 
%
Though 3D metrics are used to evaluate point cloud reconstruction quality, \gs is evaluated by 2D metrics such as PSNR~\cite{psnr} and SSIM~\cite{ssim}.
%
The views (ones not used for training) from reconstructed \gs are compared with the ground truth camera views.
%
Future work can evaluate \livo and \gsnfs using other metrics.
